Write $e_I$ for the basis blade of $\Cl(p, q, r)$ on the increasing tuple $I$ (\claimref{std-algebra.blade-basis}). \texttt{GA.blade\_mul(I, J)} starts from the concatenated word $I J$ with coefficient $1$ and repeatedly applies one of two moves to an adjacent pair of indices: if $i > j$ it swaps them and negates the coefficient; if $i = j$ it deletes both and multiplies the coefficient by $Q(e_i)$, stopping with coefficient $0$ when that is $0$. Read the word with its coefficient as the element $c\, e_{k_1} \cdots e_{k_m}$ of $\Cl(p, q, r)$. Each move leaves that element unchanged: a swap of distinct orthogonal generators because $e_i e_j = -e_j e_i$ (\claimref{std-algebra.orthogonal-anticommute}), a deletion because $e_i e_i = Q(e_i)$ (the defining relation of \claimref{obj.std-algebra.clifford-algebra}). Every move lowers the number of inversions or the length, so the procedure stops, and it stops only at an increasing word without repeats, that is at $c\, e_K$. By induction on the number of moves, \texttt{blade\_mul(I, J)} returns $(c, K)$ with $e_I e_J = c\, e_K$ in $\Cl(p, q, r)$, for every $n$. \texttt{GA.mul} extends this bilinearly.

So the linear bijection from $\Cl(p, q, r)$ to the library's space of blade tuples sending $e_I$ to $I$ carries the Clifford product to \texttt{GA.mul}. It is an isomorphism of algebras, and in particular \texttt{GA.mul} is associative with unit $()$ and satisfies $e_i^{2} = Q(e_i)$ and $e_i e_j = -e_j e_i$ for $i \ne j$, because the Clifford product does. The check is an independent cross-check of the code: it runs associativity and the relations exhaustively on basis triples for $n \le 3$ and on sampled triples for $n = 4$. The universal property (\claimref{std-algebra.existence-uniqueness}) is what makes $\Cl(p, q, r)$, and so the blade basis used here, well defined.
